% SOURCE_AUTHORITY: sga5_fr_workpass.tex, lines 5573--5608 (LF-normalized unit).
% SOURCE_SHA256: 791F4EFFC5E02832D5D77ED03518C8156D6F07E4C8238B03545DB93D883FBB28
Por fim, sejam, para $i=1,2$, $A_i$ uma $\Lambda$-álgebra, sendo $A_2$ plana, $E_1\in\Ob D_{\mathrm{ctf}}({}_{A_1}X_1)$, $F_2\in\Ob D^+({}_{A_2}X_2)$, $F_1\in\Ob D({}_{A_1}X_{1A_2})$, com $F_1$ de dimensão Tor finita e com cohomologia construtível como objeto de $D(X_{1A_2})$. Então, tem-se
\[
\uRHom_{A_1}(E_1,F_1)\in\Ob D_{\mathrm{ctf}}(X_{1A_2}),
\]
como se viu anteriormente, e dispõe-se de uma flecha canônica de $D^+(X)$ ($D^+(X,A_2)$ quando $A_2$ é comutativa),
\begin{equation}
\uRHom_{A_1}(E_1,F_1)\lten_{S,A_2}F_2
\longrightarrow
\uRHom_{A_1}(p_1^*E_1,F_1\lten_{S,A_2}F_2),
\tag{6.2.3}
\end{equation}
definida compondo a flecha de ``imagem inversa'' (III 2.1.1)
\[
p_1^*\uRHom_{A_1}(E_1,F_1)\lten_{A_2}p_2^*F_2
\longrightarrow
\uRHom_{A_1}(p_1^*E_1,p_1^*F_1)\lten_{A_2}p_2^*F_2
\]
com a flecha semelhante a 5.4.2
\[
\uRHom_{A_1}(p_1^*E_1,p_1^*F_1)\lten_{A_2}p_2^*F_2
\longrightarrow
\uRHom_{A_1}(p_1^*E_1,p_1^*F_1\lten_{A_2}p_2^*F_2)
\]
(adjunta, segundo 5.2.3, da flecha
\[
p_1^*E_1\lten\uRHom_{A_1}(p_1^*E_1,p_1^*F_1)\lten_{A_2}p_2^*F_2
\longrightarrow
p_1^*F_1\lten_{A_2}p_2^*F_2
\]
definida pela flecha de avaliação 5.3.6). De 6.2.1 deduz-se, como em (III 2.3), que 6.2.3 é um isomorfismo. Em particular, para $A_2=\Lambda$, $F_2=\Lambda_{X_2}$, 6.2.3 dá um isomorfismo
\begin{equation}
p_1^*\uRHom_{A_1}(E_1,F_1)
\xrightarrow{\sim}
\uRHom_{A_1}(p_1^*E_1,p_1^*F_1).
\tag{6.2.4}
\end{equation}
